%----------------------------------------------------------------------------------------
%    PACKAGES AND THEMES
%----------------------------------------------------------------------------------------

\documentclass[aspectratio=169,xcolor=dvipsnames]{beamer}

\AtBeginSection[]
{
    \begin{frame}
        \frametitle{Visão Geral}
        \tableofcontents[currentsection]
    \end{frame}
}

\usefonttheme{professionalfonts} % using non standard fonts for beamer
\usefonttheme{serif} % default family is serif
\usepackage{fontspec}
\setmainfont{XCharter}[
    UprightFont    = *-Roman,
    BoldFont       = *-Bold,
    ItalicFont     = *-Italic,
    BoldItalicFont = *-BoldItalic,
]

\usepackage{hyperref}
\usepackage{graphicx} % Allows including images
\usepackage{booktabs} % Allows the use of \toprule, \midrule and \bottomrule in tables
\input{./slides/SimplePlusTheme/beamerthemeSimplePlus.sty}
\input{./slides/SimplePlusTheme/beamerinnerthemeSimplePlus.sty}
\input{./slides/SimplePlusTheme/beamerfontthemeSimplePlus.sty}
\input{./slides/SimplePlusTheme/beamercolorthemeSimplePlus.sty}

\usepackage{listings}
\setbeamercovered{transparent}

\usepackage{biblatex}
\addbibresource{slides/reference.bib}
\renewcommand*{\bibfont}{\footnotesize}
% Tamanho dos slides de "Leituras Recomendadas": maior que a bibliografia
% completa (\bibfont acima, que rende \tiny via a redefinição de \footnotesize)
% para destacar as leituras curadas.
\newcommand{\leiturasize}{\small}

\usepackage[brazilian ]{babel}
\usepackage{csquotes}
\usepackage{pgfplots}
\pgfplotsset{compat=1.18}

\usepackage{emoji}
\usepackage{amsmath}
\usepackage[xcharter,bigdelims,vvarbb]{newtxmath}
% newtxmath's operators font requests the fontspec family `XCharter(0)' in OT1;
% without these substitutions math digits and operator names silently fall back
% to Computer Modern (LaTeX Font Warning: Font shape `OT1/XCharter(0)/m/n' undefined).
\DeclareFontFamily{OT1}{XCharter(0)}{}
\DeclareFontShape{OT1}{XCharter(0)}{m}{n}{<->ssub * XCharter-TLF/m/n}{}
\DeclareFontShape{OT1}{XCharter(0)}{m}{it}{<->ssub * XCharter-TLF/m/it}{}
\DeclareFontShape{OT1}{XCharter(0)}{b}{n}{<->ssub * XCharter-TLF/b/n}{}
\DeclareFontShape{OT1}{XCharter(0)}{bx}{n}{<->ssub * XCharter-TLF/b/n}{}

\usepackage{bm}
\usepackage[table]{xcolor}
\usepackage{xcolor}
\usepackage{algpseudocode}

\usepackage{cancel}
\usepackage{ulem}

\usetikzlibrary{shapes}
\usetikzlibrary{arrows}
\usetikzlibrary {arrows.meta}
\usetikzlibrary{calc,positioning}
\usepgfplotslibrary{fillbetween}
\usetikzlibrary{angles,quotes}
\usetikzlibrary{tikzmark}


\definecolor{PastelRed}{HTML}{FFC1CC}
\definecolor{PastelBlue}{HTML}{C2F0FF}
\definecolor{PastelBlue2}{HTML}{3182BD}

\definecolor{PastelPink}{HTML}{FFCBE1}
\definecolor{PastelGreen}{HTML}{D6E5BD}
\definecolor{PastelYellow}{HTML}{F9E1A8}
\definecolor{PastelBlue}{HTML}{BCD8EC}
\definecolor{PastelPurple}{HTML}{DCCCEC}
\definecolor{PastelOrange}{HTML}{FFDAB4}


\renewcommand{\footnotesize}{\tiny}

\newcommand{\mespor}{
  \ifcase\month
    \or janeiro
    \or fevereiro
    \or março
    \or abril
    \or maio
    \or junho
    \or julho
    \or agosto
    \or setembro
    \or outubro
    \or novembro
    \or dezembro
  \fi
}
% \usepackage{csquotes}
% \usepackage{minted}
% \usemintedstyle{xcode}
\definecolor{vscLightBackground}{HTML}{FFFFFF}
\definecolor{vscLightText}{HTML}{000000}
\definecolor{vscLightGreen}{HTML}{008000}        % For Comments AND Numbers
\definecolor{vscLightRed}{HTML}{A31515}          % For Strings
\definecolor{vscLightTeal}{HTML}{267F99}
\definecolor{vscBlue}{HTML}{0000FF}% For main keywords (if, for, def...)
\definecolor{vscLightPytorch}{HTML}{800080}      % For PyTorch keywords (dark magenta)

\lstdefinestyle{vscode-light}{
    language=Python,
    backgroundcolor=\color{vscLightBackground},
    basicstyle=\ttfamily\scriptsize\color{vscLightText},
    keywordstyle=\color{vscLightTeal},
    commentstyle=\color{vscLightGreen},
    stringstyle=\color{vscLightRed},
    % Style for PyTorch keywords (using keyword group 2)
    keywordstyle=[2]\color{vscLightPytorch}, 
    morekeywords={self, True, False, None},
    morekeywords=[2]{
        torch, Tensor, device, dtype, to, nn, optim, autograd, utils, data,
        Module, Linear, Conv2d, ReLU, CrossEntropyLoss, Adam, SGD,
        Dataset, DataLoader, randn, zeros, ones, arange, cat, stack,
        view, reshape, sum, mean, max, min, matmul, no_grad, backward
    },
    keywordstyle=[3]\color{vscBlue}, 
    morekeywords=[3]{
        model, loss, y, y_hat, x , optimizer
    },
    frame=single,
    tabsize=1,
    showstringspaces=false,
    breaklines=true,
    captionpos=b,
    extendedchars=true,
    numbers=left,
}

\usepackage[cache=true]{minted}
% minted v3+ (TeX Live 2024+) auto-detects the output directory.
% minted v2 (Ubuntu apt TeX Live 2023) needs it set explicitly to
% match latexmkrc's $out_dir = 'build'.
\makeatletter
\@ifpackagelater{minted}{2024/01/01}{}{%
  \def\minted@outputdir{build/}%
}
\makeatother
\setminted{
    fontsize=\scriptsize,
    style=vs,
    breaklines=true,
    numbers=left,
    numbersep=5pt,
    frame=single,
}
\providecommand{\citeauthoryear}[1]{(\citeauthor{#1},~\citeyear{#1})}

\providecommand{\thetav}{\bm{\theta}}

% vetor coluna 2D com colchetes
\providecommand{\cvec}[2]{\begin{bmatrix} #1 \\ #2 \end{bmatrix}}

\providecommand{\varvector}[1]{\mathbf{#1}}
% instance
\providecommand{\X}{\varvector{X}}
\providecommand{\mlinstance}{\varvector{X}}
\providecommand{\mlinstancei}{\varvector{X}^{(i)}}
\providecommand{\mlinstancex}[1]{\varvector{X}^{(#1)}}

% target
\providecommand{\y}{\varvector{y}}
\providecommand{\mltarget}{\varvector{y}}
\providecommand{\mltargeti}{\varvector{y}^{(i)}}
\providecommand{\mltargetx}[1]{\varvector{y}^{(#1)}}

%hats
\providecommand{\yhat}{\hat{\varvector{y}}}
\providecommand{\mltargethat}{\hat{\varvector{y}}}
\providecommand{\mltargethati}{\hat{\varvector{y}}^{(i)}}
\providecommand{\mltargethatx}[1]{\hat{\varvector{y}}^{(#1)}}

% linear reg
\providecommand{\mllinearregression}{\mltargethat=\mlinstance\bm{\theta}}

\providecommand{\tikzarrowcolor}[3]{%
  \begin{tikzpicture}[remember picture, overlay]
    % Arrow body
    \path[fill=#3] ([shift={(#1,#2)}]current page.south west) 
      rectangle ++(0.3, -0.2);
    % Arrow head
    \path[fill=#3] ([shift={(#1+0.3,#2+0.1)}]current page.south west) -- 
         ++(0, -0.4) -- 
         ++(0.2, 0.2) -- 
         cycle;
  \end{tikzpicture}%
}

\providecommand{\tikzarrow}[2]{%
    \tikzarrowcolor{#1}{#2}{Red}
}



\DeclareMathOperator*{\argmax}{argmax}
\DeclareMathOperator*{\argmin}{argmin}



%----------------------------------------------------------------------------------------
%    TITLE PAGE
%----------------------------------------------------------------------------------------

\title{Deep Learning}
\subtitle{Aplicações de CNN em Visão Computacional}

\author{Lucas Silveira Kupssinskü}

\institute
{
    Machine Learning Theory and Applications Laboratory \\
    PUCRS
}
\date{\mespor de \the\year}

%----------------------------------------------------------------------------------------
%    PRESENTATION SLIDES
%----------------------------------------------------------------------------------------

\begin{document}

\begin{frame}
    \titlepage
\end{frame}

\begin{frame}{Visão Geral}
    \tableofcontents
\end{frame}

%========================================================================
\section{VGG}
%========================================================================

\begin{frame}{VGG}
\begin{itemize}
    \item Arquitetura de CNN proposta em 2014, publicada no ICLR 2015\footfullcite{simonyan2014vgg}.
    \item Continua sendo usada em muitos cenários, especialmente como \textit{backbone} pré-treinado.
    \item No ILSVRC-2014 (\textit{ImageNet Large Scale Visual Recognition Challenge}) ficou em \textbf{1º na localização} e \textbf{2º na classificação} (vencida pelo GoogLeNet).
    \item Taxa de erro no ImageNet (ensemble multi-crop reportado no paper):
    \begin{itemize}
        \item \textit{Top-5}: $6.8\%$.
        \item \textit{Top-1}: $23.7\%$.
        \item VGG-16 \textit{single-model} fica em torno de $7.3\%$ \textit{top-5} e $24.7\%$ \textit{top-1}.
    \end{itemize}
    \item $\sim 140$M parâmetros treináveis.
\end{itemize}
\end{frame}

\begin{frame}{Definição do Problema}
\begin{itemize}
    \item ImageNet contém imagens RGB na resolução $224\times 224 \times 3$.
    \begin{itemize}
        \item As entradas são tensores $(224, 224, 3)$.
    \end{itemize}
    \item O objetivo é classificar cada imagem em uma das $1000$ classes pré-definidas.
    \begin{itemize}
        \item Função de custo: entropia cruzada com os \textit{logits}.
    \end{itemize}
    \item Precisamos transformar um tensor $(224, 224, 3)$ em um vetor $(1, 1, 1000)$.
\end{itemize}
\end{frame}

\begin{frame}{Blocos de Construção}
\begin{itemize}
    \item A VGG combina três tipos de operações:
    \begin{itemize}
        \item \textbf{Convolução + ReLU}: aumenta o número de canais e extrai características espaciais.
        \item \textbf{\textit{Pooling}}: subamostra a dimensão espacial (largura e altura).
        \item \textbf{Camadas densas}: realizam a classificação final, partindo do \textit{feature map} achatado.
    \end{itemize}
    \item Toda a parte convolucional pode ser vista como um \textbf{extrator de \textit{features}}.
    \item As camadas densas no fim convertem essas \textit{features} em probabilidades de classe.
\end{itemize}
\end{frame}

\begin{frame}{VGG16}
\begin{itemize}
    \item VGG16 é a implementação mais comum, equivalente à variante ``D'' descrita no artigo.
    \item Possui 16 camadas com parâmetros treináveis: \textcolor{BrickRed}{\textbf{13 convolucionais}} e \textcolor{PastelBlue2}{\textbf{3 densas}}.
    \item Todas as convoluções são $3\times 3$ com \textit{stride} 1 e \textit{padding} 1 (mantêm a dimensão espacial).
    \item Cada \textit{max pooling} usa janela $2\times 2$ com \textit{stride} 2 (reduz pela metade a largura e altura).
\end{itemize}
\end{frame}

\begin{frame}{Arquitetura da VGG16}
\centering
\begin{tikzpicture}[every node/.style={font=\scriptsize}, scale=0.95]
    % Helper: column index x position for the spatial pyramid
    % Use simple cuboid-like rectangles to suggest tensor volumes.
    \def\bw{0.55} % box width

    % Input
    \draw[fill=black!5] (0,0) rectangle ++(\bw,2.6);
    \node[align=center] at (\bw/2,-0.35) {$224{\times}224$\\$\times 3$};

    % Block 1: conv 64 x2, then pool
    \draw[fill=BrickRed!25] (1.0,0) rectangle ++(\bw,2.6);
    \draw[fill=BrickRed!25] (1.6,0) rectangle ++(\bw,2.6);
    \node[align=center] at (1.55,-0.35) {$224{\times}224$\\$\times 64$};

    % After pool
    \draw[fill=BrickRed!40] (2.7,0.3) rectangle ++(\bw,2.0);
    \draw[fill=BrickRed!40] (3.3,0.3) rectangle ++(\bw,2.0);
    \node[align=center] at (3.25,-0.35) {$112{\times}112$\\$\times 128$};

    % After pool
    \draw[fill=BrickRed!55] (4.4,0.55) rectangle ++(\bw,1.5);
    \draw[fill=BrickRed!55] (5.0,0.55) rectangle ++(\bw,1.5);
    \draw[fill=BrickRed!55] (5.6,0.55) rectangle ++(\bw,1.5);
    \node[align=center] at (5.25,-0.35) {$56{\times}56$\\$\times 256$};

    % After pool
    \draw[fill=BrickRed!70] (6.7,0.75) rectangle ++(\bw,1.1);
    \draw[fill=BrickRed!70] (7.3,0.75) rectangle ++(\bw,1.1);
    \draw[fill=BrickRed!70] (7.9,0.75) rectangle ++(\bw,1.1);
    \node[align=center] at (7.55,-0.35) {$28{\times}28$\\$\times 512$};

    % After pool
    \draw[fill=BrickRed!85] (9.0,0.95) rectangle ++(\bw,0.7);
    \draw[fill=BrickRed!85] (9.6,0.95) rectangle ++(\bw,0.7);
    \draw[fill=BrickRed!85] (10.2,0.95) rectangle ++(\bw,0.7);
    \node[align=center] at (9.85,-0.35) {$14{\times}14$\\$\times 512$};

    % After last pool
    \draw[fill=BrickRed!100, draw=black] (11.3,1.15) rectangle ++(\bw,0.3);
    \node[align=center] at (11.55,-0.35) {$7{\times}7$\\$\times 512$};

    % Dense layers
    \draw[fill=PastelBlue2, opacity=0.7] (12.4,1.2) rectangle ++(0.25,0.2);
    \draw[fill=PastelBlue2, opacity=0.7] (12.85,1.2) rectangle ++(0.25,0.2);
    \node[align=center] at (12.75,-0.35) {fc\\4096};

    % Output (softmax)
    \draw[fill=Goldenrod, opacity=0.7] (13.4,1.2) rectangle ++(0.2,0.2);
    \node[align=center] at (13.5,-0.35) {soft\\max};

    % Legend
    \node[anchor=west, font=\tiny] at (0,3.0) {\tikz \fill[BrickRed!50] (0,0) rectangle (0.25,0.18); ~conv $3\times 3$ + ReLU \quad
        \tikz \fill[PastelBlue2, opacity=0.7] (0,0) rectangle (0.25,0.18); ~totalmente conectada \quad
        \tikz \fill[Goldenrod, opacity=0.7] (0,0) rectangle (0.25,0.18); ~softmax};
\end{tikzpicture}

\vspace{0.3em}
\footnotesize Cada \textit{max pooling} reduz a dimensão espacial pela metade e dobra (ou estabiliza) o número de canais.
\end{frame}

\begin{frame}{Contando Parâmetros: Primeiro Bloco}
\begin{itemize}
    \item Bloco inicial: duas convoluções $3\times 3$ produzindo 64 canais.
    \begin{itemize}
        \item Cada filtro convolucional tem $K_h \cdot K_w \cdot C_{in}$ pesos, e há $C_{out}$ filtros.
        \item Primeira: kernel $3\times 3$, $C_{in}=3$ (RGB), $C_{out}=64$, portanto $3 \cdot 3 \cdot 3 \cdot 64 = 9 \cdot 192 = 1.728$ pesos.
        \item Segunda: kernel $3\times 3$, $C_{in}=64$ (saída da primeira), $C_{out}=64$, portanto $3 \cdot 3 \cdot 64 \cdot 64 = 9 \cdot 4.096 = 36.864$ pesos.
    \end{itemize}
    \item Total parcial: $1.728 + 36.864 = 38.592$ parâmetros.
    \item O \textit{max pooling} não adiciona parâmetros.
\end{itemize}
\begin{block}{Observação}
Ignoramos o \textit{bias} para simplificar a contagem. Na prática, cada filtro tem um \textit{bias} associado.
\end{block}
\end{frame}

\begin{frame}{Contando Parâmetros: Convoluções}
\begin{itemize}
    \item Crescimento por bloco convolucional:
\end{itemize}
\vspace{0.2em}
\centering
\begin{tabular}{lccc}
\toprule
Bloco & Convoluções & Parâmetros do bloco & Total acumulado \\
\midrule
1 & $2 \times$ conv $3\times 3 \times 64$ & $38.592$ & $38.592$ \\
2 & $2 \times$ conv $3\times 3 \times 128$ & $221.184$ & $259.776$ \\
3 & $3 \times$ conv $3\times 3 \times 256$ & $1.474.560$ & $1.734.336$ \\
4 & $3 \times$ conv $3\times 3 \times 512$ & $5.898.240$ & $7.632.576$ \\
5 & $3 \times$ conv $3\times 3 \times 512$ & $7.077.888$ & $14.710.464$ \\
\bottomrule
\end{tabular}
\vspace{0.6em}
\begin{itemize}
    \item Após cinco blocos convolucionais: $\sim 14.7$M parâmetros.
\end{itemize}
\end{frame}

\begin{frame}{Contando Parâmetros: Camadas Densas}
\begin{itemize}
    \item Após o último \textit{pooling}, achatamos um tensor $7\times 7\times 512 = 25.088$.
\end{itemize}
\vspace{0.2em}
\centering
\begin{tabular}{lcc}
\toprule
Camada & Cálculo & Parâmetros \\
\midrule
fc1 & $25.088 \times 4.096$ & $102.760.448$ \\
fc2 & $4.096 \times 4.096$ & $16.777.216$ \\
fc3 (saída) & $4.096 \times 1.000$ & $4.096.000$ \\
\midrule
\textbf{Subtotal denso} & & $\mathbf{123.633.664}$ \\
\bottomrule
\end{tabular}
\vspace{0.6em}
\begin{itemize}
    \item Total da rede: $14.710.464 + 123.633.664 \approx 138.3$M parâmetros.
\end{itemize}
\end{frame}

\begin{frame}{Onde Estão os Parâmetros?}
\begin{columns}
    \begin{column}{.55\linewidth}
        \begin{itemize}
            \item A maioria dos parâmetros treináveis está nas \textbf{camadas densas}.
            \begin{itemize}
                \item Convoluções: $\sim 14$M.
                \item Densas: $\sim 124$M.
            \end{itemize}
            \item Convoluções compartilham pesos, então usam ordens de magnitude menos parâmetros que camadas totalmente conectadas equivalentes.
            \item Por isso é comum trocar as densas por uma cabeça menor quando se reaproveita a VGG em outras tarefas.
        \end{itemize}
    \end{column}
    \begin{column}{.45\linewidth}
        \centering
        \begin{tikzpicture}[every node/.style={font=\scriptsize}]
            % Stacked horizontal bar for param distribution
            \def\barH{0.6}
            \def\totalW{6}
            \pgfmathsetmacro{\convW}{\totalW * 14 / 138}
            \pgfmathsetmacro{\denseW}{\totalW * 124 / 138}
            % Conv slice
            \fill[BrickRed!60] (0,0) rectangle (\convW,\barH);
            \fill[PastelBlue2!70] (\convW,0) rectangle (\totalW,\barH);
            \draw[thick] (0,0) rectangle (\totalW,\barH);
            % Labels above
            \node[anchor=south, BrickRed] at (\convW/2,\barH) {\tiny $\sim 14$M};
            \node[anchor=south, PastelBlue2] at (\convW + \denseW/2,\barH) {\tiny $\sim 124$M};
            % Legend below
            \node[anchor=north west, font=\tiny] at (0,-0.2) {\tikz \fill[BrickRed!60] (0,0) rectangle (0.25,0.18); ~Convolucionais};
            \node[anchor=north west, font=\tiny] at (0,-0.7) {\tikz \fill[PastelBlue2!70] (0,0) rectangle (0.25,0.18); ~Densas};
        \end{tikzpicture}
    \end{column}
\end{columns}
\end{frame}

%========================================================================
\section{Localização}
%========================================================================

\begin{frame}{Classificação vs Classificação + Localização}
\begin{itemize}
    \item Além de dizer ``há um gato na imagem'', queremos saber \textbf{onde} esse gato está.
    \item A saída adicional é uma \textit{bounding box} delimitando o objeto.
\end{itemize}
\vspace{0.4em}
\centering
\begin{tikzpicture}[every node/.style={font=\scriptsize}]
    % Classification box
    \draw[thick] (0,0) rectangle (3,2.4);
    \node at (1.5,2.7) {\small Classificação};
    \node[anchor=south] at (1.5,1.2) {\large \textbf{Gato}};
    \node[anchor=south] at (1.5,0.2) {imagem inteira};

    % Classification + Localization
    \draw[thick] (4.5,0) rectangle (7.5,2.4);
    \node at (6,2.7) {\small Classificação + Localização};
    \draw[thick, BrickRed] (5.2,0.4) rectangle (7.0,1.9);
    \node[BrickRed] at (6.1,2.1) {gato};
    \node[anchor=south] at (6,-0.3) {classe + caixa};
\end{tikzpicture}
\end{frame}

\begin{frame}{Bounding Box}
\begin{itemize}
    \item Para descrever uma \textit{bounding box} precisamos de 4 valores:
    \begin{itemize}
        \item $(x, y)$: coordenadas de um ponto de referência (canto superior esquerdo ou centro).
        \item $h$: altura da caixa.
        \item $w$: largura da caixa.
    \end{itemize}
    \item Podemos modelar esses 4 valores como um problema de \textbf{regressão multivariada}.
    \item Convenção comum: o canto superior esquerdo da imagem é $(0,0)$ e o canto inferior direito é $(1,1)$.
\end{itemize}
\end{frame}

\begin{frame}{Exemplo de Bounding Box}
\begin{columns}
    \begin{column}{.5\linewidth}
        \begin{itemize}
            \item Alvo $\mathbf{y} = (x, y, h, w)$.
            \item Para o exemplo ao lado:
            \begin{itemize}
                \item $(x, y) = (0.2,\, 0.1)$.
                \item $h = 0.5$.
                \item $w = 0.7$.
            \end{itemize}
            \item Logo, $\mathbf{y} = (0.2,\, 0.1,\, 0.5,\, 0.7)$.
        \end{itemize}
    \end{column}
    \begin{column}{.5\linewidth}
        \centering
        \begin{tikzpicture}[every node/.style={font=\scriptsize}]
            % Image frame (unit coordinate)
            \draw[thick, fill=black!3] (0,0) rectangle (4,4);
            \node[anchor=south west, font=\tiny] at (0,4.05) {$(0, 0)$};
            \node[anchor=north east, font=\tiny] at (4,-0.05) {$(1, 1)$};
            % Bounding box
            \draw[thick, BrickRed] (0.8,1.6) rectangle (3.6,3.6);
            \node[anchor=south, BrickRed, font=\tiny] at (0.9,3.6) {$(0.2,\, 0.1)$};
            \draw[<->, BrickRed] (3.8,1.6) -- (3.8,3.6) node[midway, right, font=\tiny] {$h=0.5$};
            \draw[<->, BrickRed] (0.8,1.4) -- (3.6,1.4) node[midway, below, font=\tiny] {$w=0.7$};
        \end{tikzpicture}
    \end{column}
\end{columns}
\end{frame}

\begin{frame}{Arquitetura para Localização}
\begin{itemize}
    \item Para regressão dos 4 valores da \textit{bounding box}, a saída precisa de \textbf{4 unidades sigmoide} (em vez de \textit{softmax} para 1000 classes).
    \begin{itemize}
        \item A sigmoide garante saídas em $[0, 1]$, compatíveis com coordenadas normalizadas.
        \item Hipótese implícita: assumimos que a caixa cabe inteiramente dentro da imagem, com $(x, y) \in [0,1]^2$ e $h, w \in [0,1]$.
        \item Caixas que estouram a borda exigem outra parametrização (por exemplo, logits sem ativação, ou regressão de \textit{offsets} em relação a uma \textit{anchor}).
    \end{itemize}
    \item A função de custo passa a ser o \textbf{erro quadrático médio} (MSE):
\end{itemize}
\begin{equation*}
\mathcal{L} = \frac{1}{2N}\sum_{i=1}^{N}\Bigl[\bigl(y_0^{(i)} - \hat{y}_0^{(i)}\bigr)^2 + \bigl(y_1^{(i)} - \hat{y}_1^{(i)}\bigr)^2 + \bigl(y_2^{(i)} - \hat{y}_2^{(i)}\bigr)^2 + \bigl(y_3^{(i)} - \hat{y}_3^{(i)}\bigr)^2\Bigr]
\end{equation*}
\begin{itemize}
    \item As convoluções continuam funcionando como extrator de \textit{features}.
\end{itemize}
\end{frame}

\begin{frame}{Reaproveitando uma Rede Pré-treinada}
\begin{itemize}
    \item Imagine que já treinamos a rede para \textbf{classificar} as fotos.
    \begin{itemize}
        \item Os parâmetros das convoluções já estão ajustados para extrair \textit{features} úteis.
    \end{itemize}
    \item Estratégia comum: \textbf{congelar} os pesos da convolução e treinar apenas a cabeça densa para a nova tarefa\footfullcite{sermanet2013overfeat}.
    \begin{itemize}
        \item Reduz drasticamente o tempo de treinamento.
        \item Funciona bem quando o conjunto de dados novo é pequeno.
    \end{itemize}
    \item Alternativa: \textit{fine-tuning} da rede inteira com \textit{learning rate} baixa.
\end{itemize}
\end{frame}

\begin{frame}{Métricas de Localização: \textit{Intersection over Union}}
\begin{columns}
    \begin{column}{.55\linewidth}
        \begin{itemize}
            \item Computa a interseção entre a \textit{bounding box} prevista e o \textit{ground truth}, sobre a união das duas áreas.
            \begin{equation*}
                IoU = \frac{A \cap B}{A \cup B}
            \end{equation*}
            \item Valor ideal: $IoU = 1$ (sobreposição perfeita).
            \item Ignora o tamanho absoluto das caixas, só importa a sobreposição relativa.
        \end{itemize}
    \end{column}
    \begin{column}{.45\linewidth}
        \centering
        \begin{tikzpicture}
            % Ground truth
            \draw[thick, BrickRed, fill=BrickRed, fill opacity=0.25] (0,0.6) rectangle (2.6,2.6);
            \node[BrickRed, font=\scriptsize] at (0.5,2.85) {ground truth $A$};
            % Prediction
            \draw[thick, PastelBlue2, fill=PastelBlue2, fill opacity=0.25] (1.4,0) rectangle (3.6,2.0);
            \node[PastelBlue2, font=\scriptsize] at (3.0,-0.3) {predição $B$};
            % Intersection emphasis
            \draw[ultra thick, dashed, black] (1.4,0.6) rectangle (2.6,2.0);
            \node[font=\scriptsize] at (2.0,1.3) {$A \cap B$};
        \end{tikzpicture}
    \end{column}
\end{columns}
\end{frame}

\begin{frame}{Métricas de Localização: \textit{Dice Score}}
\begin{columns}
    \begin{column}{.55\linewidth}
        \begin{itemize}
            \item Computa o dobro da interseção sobre a soma das áreas.
            \begin{equation*}
                Dice = \frac{2 \cdot (A \cap B)}{Area(A) + Area(B)}
            \end{equation*}
            \item Valor ideal: $Dice = 1$.
            \item Penaliza menos diferenças de tamanho que o IoU.
        \end{itemize}
    \end{column}
    \begin{column}{.45\linewidth}
        \centering
        \begin{tikzpicture}
            \draw[thick, BrickRed, fill=BrickRed, fill opacity=0.25] (0,0.6) rectangle (2.6,2.6);
            \node[BrickRed, font=\scriptsize] at (0.5,2.85) {$A$};
            \draw[thick, PastelBlue2, fill=PastelBlue2, fill opacity=0.25] (1.4,0) rectangle (3.6,2.0);
            \node[PastelBlue2, font=\scriptsize] at (3.4,-0.3) {$B$};
            \draw[ultra thick, dashed, black] (1.4,0.6) rectangle (2.6,2.0);
        \end{tikzpicture}
    \end{column}
\end{columns}
\end{frame}

\begin{frame}{IoU vs Dice: Exemplo Numérico}
\begin{itemize}
    \item Suponha duas caixas com:
    \begin{itemize}
        \item $Area(A) = 10$.
        \item $Area(B) = 10$.
        \item $A \cap B = 5$.
        \item $A \cup B = Area(A) + Area(B) - A \cap B = 15$.
    \end{itemize}
\end{itemize}
\vspace{0.2em}
\begin{columns}
    \begin{column}{.5\linewidth}
        \begin{equation*}
            IoU = \frac{5}{15} \approx 0.333
        \end{equation*}
    \end{column}
    \begin{column}{.5\linewidth}
        \begin{equation*}
            Dice = \frac{2 \cdot 5}{10 + 10} = 0.5
        \end{equation*}
    \end{column}
\end{columns}
\vspace{0.6em}
\begin{itemize}
    \item Ambas as métricas medem a sobreposição entre caixas.
    \item Para a mesma sobreposição, o IoU produz valores menores que o \textit{Dice}.
    \begin{itemize}
        \item Pode ser interpretado como uma penalização ``mais quadrática'' do erro.
    \end{itemize}
\end{itemize}
\end{frame}

\begin{frame}{Localização: Resumo}
\begin{itemize}
    \item A \textit{bounding box} é modelada como quatro saídas de regressão multivariada.
    \item Função de custo: MSE.
    \item Ativação na saída: sigmoide (4 unidades).
    \item Podemos:
    \begin{itemize}
        \item Reaproveitar pesos de convolução já treinados.
        \item Treinar apenas a nova cabeça densa.
        \item Ou treinar a rede inteira (\textit{fine-tuning}).
    \end{itemize}
    \item Avaliamos com IoU ou \textit{Dice Score}, comparando com o \textit{ground truth}.
\end{itemize}
\end{frame}

%========================================================================
\section{Classificação e Localização}
%========================================================================

\begin{frame}{Múltiplas Tarefas Simultâneas}
\begin{itemize}
    \item Podemos resolver classificação e localização \textbf{ao mesmo tempo}.
    \item A parte convolucional serve como \textbf{extrator de \textit{features}} compartilhado.
    \item Adicionamos duas cabeças densas independentes:
    \begin{itemize}
        \item Cabeça de classificação: \textit{softmax} sobre as 1000 classes.
        \item Cabeça de localização: 4 unidades sigmoide para a \textit{bounding box}.
    \end{itemize}
    \item Cada cabeça contribui com seu próprio termo na função de custo:
\end{itemize}
\begin{equation*}
\mathcal{L} = \mathcal{L}_{\text{classif.}} + \lambda\,\mathcal{L}_{\text{loc.}}
\end{equation*}
\end{frame}

\begin{frame}{Backbone como Extrator de Features}
\centering
\begin{tikzpicture}[every node/.style={font=\scriptsize}]
    % Image
    \draw[thick, fill=black!5] (0,0) rectangle (1.2,2.0);
    \node[anchor=north] at (0.6,-0.05) {Imagem};
    % Conv backbone
    \draw[fill=BrickRed!40, rounded corners=2pt] (1.8,0.4) rectangle (5.4,1.6);
    \node[align=center] at (3.6,1.0) {\small Convoluções\\\scriptsize \textit{backbone} VGG};
    \draw[->, thick] (1.2,1.0) -- (1.75,1.0);
    % Output features arrow (from backbone to branch)
    \draw[thick] (5.4,1.0) -- (6.2,1.0);
    % Two heads
    \draw[fill=PastelBlue2!50, rounded corners=2pt] (6.4,1.5) rectangle (8.8,2.4);
    \node[align=center] at (7.6,2.0) {\small Cabeça densa\\\scriptsize classificação};
    \draw[fill=PastelBlue2!50, rounded corners=2pt] (6.4,-0.4) rectangle (8.8,0.5);
    \node[align=center] at (7.6,0.05) {\small Cabeça densa\\\scriptsize localização};
    \draw[->, thick] (6.2,1.0) -- (6.2,1.95) -- (6.4,1.95);
    \draw[->, thick] (6.2,1.0) -- (6.2,0.05) -- (6.4,0.05);
    % Outputs
    \draw[->, thick] (8.8,2.0) -- (9.6,2.0);
    \node[anchor=west, align=left] at (9.6,2.0) {\small classe\\\scriptsize \textit{softmax}};
    \draw[->, thick] (8.8,0.05) -- (9.6,0.05);
    \node[anchor=west, align=left] at (9.6,0.05) {\small caixa\\\scriptsize $(x,y,h,w)$};
\end{tikzpicture}

\vspace{0.4em}
\begin{itemize}
    \item Trocando a cabeça e mantendo o \textit{backbone}, podemos resolver muitas tarefas distintas a partir das mesmas \textit{features}.
\end{itemize}
\end{frame}

%========================================================================
\section{Segmentação Semântica}
%========================================================================

\begin{frame}{Segmentação Semântica}
\begin{itemize}
    \item É basicamente uma classificação em \textbf{nível de pixel}.
    \item Cada pixel da imagem recebe uma probabilidade de pertencer a cada classe.
    \item Aplicações típicas: direção autônoma, imagens médicas, sensoriamento remoto.
\end{itemize}
\vspace{0.5em}
\centering
\begin{tikzpicture}[every node/.style={font=\scriptsize}]
    \node[inner sep=0] (in)
      {\includegraphics[width=4.7cm]{slides/img/02-cnn/segmentation_example.jpg}};
    \node[above=2pt of in] {Imagem de entrada};
    \node[inner sep=0, right=1.3cm of in] (out)
      {\includegraphics[width=4.7cm]{slides/img/02-cnn/segmentation_example_mask.png}};
    \node[above=2pt of out] {Máscara por classe};
    \draw[->, very thick] (in.east) -- (out.west);
\end{tikzpicture}\\[4pt]
{\tiny Cada pixel recebe uma classe (\textit{notebook}, \textit{copo}, \textit{mesa}, fundo). Exemplo CC0 de Martin Thoma, \textit{Wikimedia Commons}.}
\end{frame}

\begin{frame}{Primeira Ideia (Ingênua)}
\begin{itemize}
    \item Separar a imagem em \textit{patches}.
    \item Para cada \textit{patch}, classificar o pixel central com uma CNN.
    \item Repetir para todos os pixels da imagem.
    \item Problemas:
    \begin{itemize}
        \item Extremamente \textbf{ineficiente}: requer milhares de \textit{forwards} por imagem.
        \item Muito redundante, pois \textit{patches} vizinhos compartilham a maior parte do conteúdo.
    \end{itemize}
    \item Precisamos de uma forma de classificar todos os pixels em um único \textit{forward}.
\end{itemize}
\end{frame}

\begin{frame}{Arquitetura Encoder-Decoder}
\begin{itemize}
    \item A CNN deve classificar todos os pixels em um único \textit{forward}\footfullcite{noh2015deconv}.
    \item \textbf{Encoder}: reduz a dimensão espacial via convoluções e \textit{pooling}.
    \item \textbf{Decoder}: aumenta a dimensão espacial via \textit{unpooling} e convoluções transpostas.
\end{itemize}
\vspace{0.2em}
\centering
\begin{tikzpicture}[every node/.style={font=\tiny}]
    \def\bw{0.35}

    % Encoder pyramid (decreasing)
    \draw[fill=BrickRed!30] (0,0.0) rectangle ++(\bw,1.8);
    \draw[fill=BrickRed!40] (0.6,0.2) rectangle ++(\bw,1.4);
    \draw[fill=BrickRed!55] (1.2,0.4) rectangle ++(\bw,1.0);
    \draw[fill=BrickRed!70] (1.8,0.55) rectangle ++(\bw,0.7);
    \draw[fill=BrickRed!85] (2.4,0.7) rectangle ++(\bw,0.4);

    % Bottleneck
    \draw[fill=black!50] (3.1,0.8) rectangle ++(0.25,0.2);
    \node at (3.22,1.25) {bottleneck};

    % Decoder (increasing)
    \draw[fill=PastelBlue2, opacity=0.35] (3.6,0.7) rectangle ++(\bw,0.4);
    \draw[fill=PastelBlue2, opacity=0.5] (4.2,0.55) rectangle ++(\bw,0.7);
    \draw[fill=PastelBlue2, opacity=0.6] (4.8,0.4) rectangle ++(\bw,1.0);
    \draw[fill=PastelBlue2, opacity=0.75] (5.4,0.2) rectangle ++(\bw,1.4);
    \draw[fill=PastelBlue2, opacity=0.9] (6.0,0.0) rectangle ++(\bw,1.8);

    % Labels
    \node[align=center] at (1.2,-0.3) {\scriptsize Encoder (VGG-16)};
    \node[align=center] at (4.8,-0.3) {\scriptsize Decoder (conv transposta)};

    % Arrows
    \draw[->, thick, dashed] (-0.2,2.0) -- (3.0,2.0) node[midway, above] {redução espacial};
    \draw[->, thick, dashed] (3.4,2.0) -- (6.6,2.0) node[midway, above] {expansão espacial};
\end{tikzpicture}

\vspace{0.2em}
\begin{itemize}
    \item Saída: máscara com a mesma resolução da entrada, com uma classe por pixel.
\end{itemize}
\end{frame}

%========================================================================
\section{Detecção de Objetos}
%========================================================================

\begin{frame}{Detecção de Objetos}
\begin{itemize}
    \item Tarefa muito parecida com classificação $+$ localização.
    \item Diferença: a imagem pode conter \textbf{vários objetos} simultaneamente.
    \begin{itemize}
        \item Cada objeto detectado tem sua própria classe e \textit{bounding box}.
    \end{itemize}
    \item Desafio principal: \textbf{não sabemos de antemão quantos objetos} aparecem em cada imagem.
\end{itemize}
\vspace{0.1em}
\centering
\resizebox{5.4cm}{!}{%
\begin{tikzpicture}[every node/.style={font=\scriptsize\bfseries}]
    \node[inner sep=0, anchor=south west] (img) at (0,0)
      {\includegraphics[width=6.5cm]{slides/img/02-cnn/det_cat_dog.jpg}};
    % cachorro (esquerda)
    \draw[BrickRed, very thick] (0.20,2.68) rectangle (3.59,4.91);
    \node[BrickRed, anchor=south west, fill=white, inner sep=1pt] at (0.20,4.91) {cachorro};
    % gato (direita)
    \draw[PastelBlue2, very thick] (3.05,1.26) rectangle (5.42,3.23);
    \node[PastelBlue2, anchor=south west, fill=white, inner sep=1pt] at (3.05,3.23) {gato};
\end{tikzpicture}}\\[3pt]
{\tiny Cada objeto recebe classe $+$ \textit{bounding box}. Foto: \textit{Socks e Buddy} (Casa Branca), domínio público.}
\end{frame}

\begin{frame}{Famílias de Detectores: Two-Stage}
\begin{itemize}
    \item \textbf{R-CNN}\footfullcite{girshick2014rcnn}: gera $\sim 2.000$ \textit{region proposals} (via \textit{selective search}), recorta cada uma e roda a CNN separadamente. Lento, mas estabeleceu a base do paradigma.
    \item \textbf{Fast R-CNN}: a CNN processa a imagem inteira uma só vez; \textit{ROI pooling} extrai \textit{features} de cada região a partir do mesmo \textit{feature map}.
    \item \textbf{Faster R-CNN}\footfullcite{ren2015fasterrcnn}: substitui \textit{selective search} por uma \textbf{Region Proposal Network} (RPN), também convolucional. Treino \textit{end-to-end}.
    \item Característica comum, primeiro propõe regiões, depois classifica cada uma. Tipicamente mais precisos, porém mais lentos.
\end{itemize}
\end{frame}

\begin{frame}{Famílias de Detectores: One-Stage e Transformers}
\begin{itemize}
    \item \textbf{YOLO}\footfullcite{redmon2016yolo} (\textit{You Only Look Once}): divide a imagem em uma grade $S\times S$ e prediz, em \textbf{uma única passada}, a classe e a \textit{bounding box} para cada célula.
    \begin{itemize}
        \item Muito mais rápido, opera em tempo real, ao custo de menor precisão em objetos pequenos.
        \item Variantes posteriores (YOLOv3-v8) recuperam grande parte dessa diferença de precisão.
    \end{itemize}
    \item \textbf{SSD} (\textit{Single Shot MultiBox Detector}): também \textit{one-stage}, prediz caixas em vários níveis de resolução simultaneamente.
    \item \textbf{DETR}\footfullcite{carion2020detr}: formula detecção como predição direta de um conjunto de caixas usando um \textbf{Transformer}.
    \begin{itemize}
        \item Elimina \textit{anchors} e \textit{non-maximum suppression}.
        \item Treina com \textit{Hungarian matching} entre predições e \textit{ground truth}.
    \end{itemize}
\end{itemize}
\end{frame}

%========================================================================
\section{Segmentação de Instâncias}
%========================================================================

\begin{frame}{Segmentação de Instâncias}
\begin{itemize}
    \item Muito parecida com detecção de objetos:
    \begin{itemize}
        \item Cada objeto recebe um valor por pixel, e não apenas uma \textit{bounding box}.
        \item Instâncias diferentes da mesma classe são \textbf{diferenciadas} entre si.
    \end{itemize}
    \item Diferença em relação à segmentação semântica: dois gatos lado a lado aparecem como duas máscaras distintas, não como uma só.
\end{itemize}
\end{frame}

\begin{frame}{Segmentação Semântica vs Instâncias}
\centering
\begin{tikzpicture}[every node/.style={font=\scriptsize}]
    \node[inner sep=0] (sem)
      {\includegraphics[width=5.5cm]{slides/img/02-cnn/seg_semantic.png}};
    \node[above=2pt of sem] {\textbf{Segmentação Semântica}};
    \node[inner sep=0, right=0.8cm of sem] (inst)
      {\includegraphics[width=5.5cm]{slides/img/02-cnn/seg_instance.png}};
    \node[above=2pt of inst] {\textbf{Segmentação de Instâncias}};
\end{tikzpicture}\\[4pt]
{\small Na \textbf{semântica}, todos os objetos da classe ``balão'' recebem a \textbf{mesma} cor;
na de \textbf{instâncias}, cada balão recebe uma cor \textbf{distinta}.}\\[2pt]
{\tiny Figura: \textit{Mask R-CNN}, \textit{Wikimedia Commons} (CC BY-SA 4.0).}
\end{frame}

%========================================================================
\section{Referências}
%========================================================================

\begin{frame}{Leituras Recomendadas}
\begin{itemize}
    \item Sugere-se \textbf{fortemente} a leitura do Capítulo 10 de \textit{Understanding Deep Learning}~\cite{prince2023understanding}.
    \begin{itemize}
        \item Disponível em \url{https://udlbook.github.io/udlbook/}.
    \end{itemize}
    \item Os artigos abaixo são acessíveis e fornecem a visão original dos autores sobre as ideias apresentadas nessa aula:
    \begin{itemize}
        \item \textit{Very deep convolutional networks for large-scale image recognition}~\cite{simonyan2014vgg}.
        \item \textit{OverFeat: Integrated recognition, localization and detection using convolutional networks}~\cite{sermanet2013overfeat}.
        \item \textit{Learning deconvolution network for semantic segmentation}~\cite{noh2015deconv}.
    \end{itemize}
\end{itemize}
\end{frame}

\begin{frame}{Referências}
\printbibliography
\end{frame}

\end{document}
